% 逻辑回归实验：PPT / 手册公式源
% PPT 使用 matplotlib mathtext 渲染为透明 PNG；
% 手册使用 Word 原生公式（OMML），可在 Word 中双击编辑。
% 编译示例：pdflatex formulas.tex

\documentclass[11pt]{article}
\usepackage{amsmath,amssymb}
\usepackage[margin=2cm]{geometry}
\begin{document}

\section*{线性分数与 sigmoid}
\begin{align}
z &= \mathbf{w}^{\top}\mathbf{x} + b \\
p &= \sigma(z) = \frac{1}{1+e^{-z}}
\end{align}

\section*{伯努利似然与 BCE}
\begin{align}
P(y=1\mid x) &= p, \qquad P(y=0\mid x) = 1-p \\
L &= \prod_i p_i^{y_i}(1-p_i)^{1-y_i} \\
\mathrm{BCE} &= -\left[y\log p + (1-y)\log(1-p)\right]
\end{align}

\section*{正则化}
\begin{align}
R_{L1}(w) &= \lambda\sum_j |w_j| \\
R_{L2}(w) &= \lambda\sum_j w_j^2 \\
L_{\mathrm{total}} &= \mathrm{BCE} + \lambda R(w)
\end{align}

\section*{评估指标}
\begin{align}
\mathrm{Accuracy} &= \frac{TP+TN}{N} \\
\mathrm{Precision} &= \frac{TP}{TP+FP} \\
\mathrm{Recall} &= \frac{TP}{TP+FN} \\
F_1 &= \frac{2PR}{P+R}
\end{align}

\end{document}
